% =====================================================================
% Supplementary Information for Paper 1 v64
% =====================================================================
\documentclass[preprint,11pt,authoryear]{elsarticle}
\usepackage{amsmath,amssymb,amsthm}
\usepackage{booktabs}
\usepackage{array}
\usepackage{tabularx}
\usepackage{graphicx}
\graphicspath{{./}{figs/}{paper1_assessment_separation_v64_src/figs/}}
\usepackage{enumitem}
\usepackage[colorlinks=true,urlcolor=blue,linkcolor=blue,citecolor=blue]{hyperref}
\emergencystretch=3em
\setlength{\tolerance}{2000}

\renewcommand{\thesection}{S\arabic{section}}
\renewcommand{\thetable}{S\arabic{table}}
\renewcommand{\thefigure}{S\arabic{figure}}
\renewcommand{\theequation}{S\arabic{equation}}
\newtheorem{theorem}{Theorem}
\renewcommand{\thetheorem}{S\arabic{theorem}}
\newtheorem{lemma}{Lemma}
\renewcommand{\thelemma}{S\arabic{lemma}}
\newtheorem{proposition}{Proposition}
\renewcommand{\theproposition}{S\arabic{proposition}}
\theoremstyle{remark}
\newtheorem{remark}{Remark}
\renewcommand{\theremark}{S\arabic{remark}}

\newcommand{\Etyp}{E_{\mathrm{typ}}}
\newcommand{\Eend}{E_{\mathrm{end}}}
\newcommand{\Ephys}{E_{\mathrm{tube,phys}}}
\newcommand{\Vtyp}{\mathcal{V}_{\mathrm{typ}}}
\newcommand{\Vweak}{\mathcal{V}_{\mathrm{weak}}}
\newcommand{\Vphys}{\mathcal{V}_{\mathrm{phys}}}
\newcommand{\Vend}{\mathcal{V}_{\mathrm{end}}}
\newcommand{\Wp}{W_+}
\newcommand{\FP}{\mathrm{FP}_{\mathrm{agg}}}
\newcommand{\Tube}{\mathrm{Tube}}
\newcommand{\Succ}{\mathrm{Succ}}
\newcommand{\End}{\mathrm{End}}
\newcommand{\conv}{\operatorname{conv}}
\newcommand{\Rp}{\mathbb{R}^n_+}
\newcommand{\Q}{\mathcal{Q}}
\newcommand{\act}[1]{\textsc{#1}}

\journal{Journal of Multi-Criteria Decision Analysis}

\begin{document}
\begin{frontmatter}
\title{Supplementary Information for:\\ Every weighting licenses a plan, no plan satisfies the floors: a quantifier-order separation between compensatory and noncompensatory transition assessment}
\author[1]{Amin Abaee}
\ead{amin_abaee@ut.ac.ir}
\affiliation[1]{organization={Independent Researcher}, city={Tehran}, country={Iran}}
\begin{abstract}
This Supplementary Information accompanies the main article. It contains the remaining proofs (S2), the substitutability-spectrum extension in full (S3), the common-shock variant in full (S4), the resource-transition benchmark in full (S5), practitioner material (S6), the verification record (S7), positioning notes per literature (S8) and the framework extensions deposited at partial status (S9). Main-text results are cited as ``main-text Theorem~2'' etc.; results stated here are numbered S1, S2, \dots. Every number is an exact rational (or a certified rational bracket of a transcendental quantity) re-derived by the deposited verification code. Nothing in this document is an empirical claim, and no main-text proof depends on S3--S9.
\end{abstract}
\end{frontmatter}

\section{Abbreviations and the result map}

\begin{table}[!htbp]
\centering\small
\caption{Abbreviations.}
\begin{tabular}{@{}ll@{}}
\toprule
CES & constant elasticity of substitution \\
DFO & Fisheries and Oceans Canada \\
LPI & Living Planet Index \\
LRP & limit reference point \\
MCDA & multi-criteria decision analysis \\
MSY & maximum sustainable yield \\
NCAM & Northern Cod Assessment Model \\
SI & Supplementary Information \\
\bottomrule
\end{tabular}
\end{table}

\begin{table}[!htbp]
\centering\small
\caption{Map between the main-text result numbers (v64) and the numbering of the previous editions (v62 and v63, which share one numbering), for readers of the deposited verification code and the Lean module, whose comments use the v62 numbers.}
\begin{tabularx}{\textwidth}{@{}>{\raggedright\arraybackslash}p{0.42\textwidth}>{\raggedright\arraybackslash}p{0.24\textwidth}X@{}}
\toprule
Main text (v64) & Previous editions (v62/v63) & Verifier / Lean reference \\
\midrule
Lemma 1 (quantifier inclusion) & Remark 1 & -- \\
Lemma 2 (full cone lossless) & Remark 2 & Lean: full-cone identity \\
Proposition 1 (hierarchy, localisation) & Proposition 3 & grid check 10, 19; Lean \\
Proposition 2 (weight-family monotonicity) & Proposition 4 & -- \\
Theorem 1 (finite-menu geometry) & Theorem 6 & -- \\
Theorem 2 (witnessed separation) & Theorem 5 & grid checks 5--9, 11--22; Lean \\
Proposition 3 (persistence, erasure) & Remark 7, Theorem 8 & grid checks 23--25 \\
Theorem 3 (blend collapse) & Theorem 9 & Lean \\
Proposition 4 (time-sharing) & Proposition 10 & Lean \\
Proposition 5 (rescue threshold) & Proposition 11 & grid check 22 \\
Theorem 4 (substitutability ladder) & Lemma A, Lemma B, Theorems S1, S2 & $\sigma$-verifier (57 checks); Lean \\
Proposition 6 (common-shock verdict) & Theorems M1, M2, M3 & common-shock verifier (77 checks) \\
S3.2 Theorem S3 (off-diagonal criterion) & Theorem G & closure verifier (50 checks) \\
S4 Theorems S6--S7 & Theorems M4, CT & common-shock verifier \\
\bottomrule
\end{tabularx}
\end{table}

\section{Remaining proofs}\label{si:proofs}

\subsection{Persistence under hold-prefix extension (main-text Proposition 3(i))}

\begin{proof}
(Hierarchy.) Backward induction on stages. Base: at the final stage the terminal sets satisfy $G\subseteq G^w\subseteq G^{\mathrm{phys}}$ by main-text Lemma~2 and the definition of the terminal sets. Step: the accepted set of a stage is the set of states from which some action keeps the tube in the safe set and the successor in the next-stage accepted set; applying main-text Proposition~1(i) to the stage's action sets and the induction hypothesis to the next-stage accepted sets preserves the three-way inclusion at the stage level.
(Persistence of strictness.) With \act{hold} the unique prefix action, a prefix state is accepted exactly when it lies in the prefix safe set and in the next-stage accepted set; the stage-0 regions are therefore the witness regions of main-text Theorem~2 pulled back through the (identity) holds, and the strictness witnesses ($I\neq\varnothing$ and the point $(\tfrac12,\tfrac1{10},\tfrac1{10})$) persist.
(Conditionality.) Each of the three mechanisms listed in the main text -- a common typed-safe continuation for every aggregate-feasible state, coincident weak and typed terminal sets on the reachable subset, or actions coupled across stages -- is sufficient to erase the gap in a multi-stage setting, so no unconditional generalisation is asserted.
\end{proof}

\subsection{The erasure datum (main-text Proposition 3(ii))}\label{si:erasure}

\emph{Datum.} Two capital forms, floors at zero, the bridge stock absent ($x\equiv0$); states are pairs $s=(s_1,s_2)$. Stage 2 (final interval): safe set $S_0^{(2)}=\mathbb R^2$ (no path constraint), terminal sets $G^{(2)}=G^{w,(2)}=G=\{s\ge0\}$ -- coincident for every weight -- and a single action \act{repair} with tube $\{s\}$ and successor $(1,1)\in G$. Stage 1: safe set $S_0=\{s\ge0\}$; two actions with constant tubes, safe at every state of $S_0$: A1 with successor $(s_1-1,s_2+1)$ and A2 with successor $(s_1+1,s_2-1)$; the stage-1 terminal architecture is the stage-2 accepted sets. Initial state $z^*=(\tfrac25,\tfrac25)$.

\emph{In the two-stage datum the gap is erased.} \act{repair} is typed-admissible -- hence $w$-admissible for every weight -- from every state, so $\mathcal V^{(2)}_{\mathrm{typ}}=\mathcal V^{(2)}_w=\mathbb R^2$ for every $w$. Backward induction then admits both A1 and A2 at $z^*$ (safe tubes; successors in $\mathbb R^2$), so $z^*$ is typed-accepted at stage 0, and weak acceptance follows.

\emph{In the single-interval truncation the gap is present.} With the final interval collapsed to its terminal sets $G=\{s\ge0\}$ and $G^w=\{w\cdot s\ge0\}$, both successors leave $G$ ($s_1-1=-\tfrac35<0$ and $s_2-1=-\tfrac35<0$), so $z^*$ is typed-rejected. For the weak doctrine, with $r=w_2/w_1$ on the interior of the cone: A1 is $w$-admissible exactly when $(s_1-1)+r(s_2+1)\ge0$, i.e.\ $r\ge\tfrac{1-s_1}{s_2+1}=\tfrac37$; A2 exactly when $(s_1+1)+r(s_2-1)\ge0$, i.e.\ $r\le\tfrac{s_1+1}{1-s_2}=\tfrac73$; the two intervals cover $[0,\infty]$ because $\tfrac37<\tfrac73$, and the boundary weights are covered directly. Hence $z^*\in\Vweak$ in the truncation and the gap contains $z^*$. \qed

\subsection{The substitutability ladder (main-text Theorem 4)}\label{si:ladderproofs}

Notation as in main-text Section~4.5: $\lambda_i=1+s_i$, normalised weights $w$, $M_\theta(\lambda;w)=(\sum_iw_i\lambda_i^\theta)^{1/\theta}$, $\sigma=1/(1-\theta)$.

\begin{lemma}[the handshake]\label{lem:handshake}
With weights normalised, $\sum_iw_i\lambda_i\ge1\iff w\cdot s\ge0$; the $\theta=1$ member therefore reproduces the compensatory operator of the main text exactly -- the same acceptance set at every state and weight, and at the witness point the same per-weight plan assignment ($\rho_1=2/3$, $\rho_2=3/2$).
\end{lemma}
\begin{proof}
$\sum_iw_i\lambda_i=\sum_iw_i+\sum_iw_is_i=1+w\cdot s$.
\end{proof}

\begin{lemma}[collapse convention]\label{lem:collapse}
For every $\theta<1$ member, any tube value $\lambda_i\le0$ rejects the plan: a coordinate at or below the collapse level cannot be compensated by surpluses elsewhere.
\end{lemma}
\noindent The geometric mean is $0$ there, CES with $\theta<0$ is undefined there, and for $\theta\in(0,1)$ the convention is the family's defining property. The $\theta=1$ member alone compensates across collapse: at the witness point \act{fast}'s trough has $\lambda=(1/5,11/5)$, and the linear member certifies it at equal weight ($6/5\ge1$) while every $\theta<1$ member with a collapsed coordinate elsewhere rejects (e.g.\ at $(1/2,1,1)$, where the trough has $\lambda_1=0$).

\paragraph{The master equation} On the diagonal $z=(x<1,s,s)$ with $s\in(1,2)$, \act{fast}'s worst-case trough is $\lambda=(s-1,s+1)$ and \act{slow}'s is $(s+1,s-1)$; \act{staged} is inadmissible ($x<1$). Write $p=w_1$ for the weight share, $A=(s-1)^\theta$, $B=(s+1)^\theta$. For $\theta>0$, \act{fast} is accepted at $p$ iff $pA+(1-p)B\ge1$ and \act{slow} iff $pB+(1-p)A\ge1$; both are affine in $p$ and sum to $A+B$. \act{fast}'s acceptance set is a downward interval $[0,p_F]$ (its condition is decreasing in $p$ since $A<B$) and \act{slow}'s an upward interval $[p_S,1]$, with $p_S=1-p_F$ by symmetry; the cover $[0,p_F]\cup[p_S,1]=[0,1]$ holds iff $p_F\ge\tfrac12$, i.e.\ iff the condition holds at $p=\tfrac12$: $(A+B)/2\ge1$. For $\theta<0$ the inequality $M_\theta\ge1$ reads $\sum w_i\lambda_i^\theta\le1$ and the same argument gives $(A+B)/2\le1$. For $\theta=0$ the condition is $(s-1)^p(s+1)^{1-p}\ge1$ with the cover at $p=\tfrac12$: $s^2-1\ge1$. This is the averaging argument: the two plan conditions sum to a weight-independent quantity, so the equal weight is always the binding one.

\begin{proof}[Proof of main-text Theorem 4]
(i) \emph{Nesting.} For fixed positive $\lambda$ and weights, $M_\theta$ is nondecreasing in $\theta$ (the power-mean inequality); the collapse convention preserves this, since acceptance at any $\theta'<1$ forces $\lambda>0$ on the serving plan's tube, where the family is defined. Per plan and weight, hence per weight, hence for the protocol: acceptance at $\theta'$ implies acceptance at every $\theta>\theta'$, so the accepted rungs form a prefix and $\sigma^*(z)$ is well defined. \emph{Finiteness} on gap states with both coordinates in $(1,2)$: at the equal weight, \act{fast}'s condition contains the term $(s_1-1)^\theta$ and \act{slow}'s the term $(s_2-1)^\theta$, both of which tend to $\infty$ as $\theta\to-\infty$ because $s_i-1<1$; so for some finite $\theta$ neither plan serves the equal weight and the state is rejected (on the diagonal this is $(s-1)^\theta\to\infty$). \emph{Positivity}: $s_1+s_2\ge2$ accepts at $\theta=1$. \emph{Only-linear states}: if $\min(s_1,s_2)\le1$, say $s_1\le1$, then \act{fast}'s trough has $\lambda_1=s_1-1\le0$ and is rejected at every $\theta<1$ by Lemma~\ref{lem:collapse}; \act{slow} alone must then cover every weight, including the weight concentrated on coordinate $2$, where its condition is $s_2-1\ge1$, i.e.\ $s_2\ge2$, excluded on a gap state. So only the linear member certifies and $\sigma^*=\infty$.
(ii) Immediate from (i): every gap state is linear-accepted, the accepted rungs form a prefix, and finiteness (or only-linearity) supplies a positive elasticity at which the state is rejected.
(iii) Evaluate the master comparison at $s=2$: for $\theta\in(0,1)$ the accepting side is $\ge2$ and $(s-1)^\theta+(s+1)^\theta=1+3^\theta>2$ strictly; for $\theta<0$ the accepting side is $\le2$ and $1+3^\theta<2$ strictly; for $\theta=0$, $4>2$. All strictly accepting at $s=2$, so by continuity some $\varepsilon>0$ has every $s\in(2-\varepsilon,2)$ accepted at that rung, and these are gap states (typed acceptance needs $s\ge2$); the acceptance condition is monotone in $s$ on $(1,2)$ (both $A$ and $B$ move the accepting way), so the accepted set is exactly $[s^*(\theta),2)$ with $s^*(\theta)<2$. The closed forms of main-text Table~4 follow by direct evaluation ($\theta=1/2$: $\sqrt{s-1}+\sqrt{s+1}\ge2$, equality at $s=5/4$; $\theta=0$: $s^2\ge2$; $\theta=-1$: $(s^2-1)/s\ge1$; $\theta=-2$: $(s+1)^2+(s-1)^2\le2(s^2-1)^2\iff s^2\ge3$; $\theta=-m$: $(s-1)^{-m}+(s+1)^{-m}\le2$), and strict monotonicity of $s^*$ in depth is the nesting of (i) applied to the critical floors. The Leontief member's acceptance is weight-independent ($\min_i\lambda_i\ge1$ does not involve $w$), so its protocol collapses into the common-plan criterion -- the typed operator; hence $\inf_z\sigma^*(z)=0$ and only $\sigma=0$ is uniformly safe.
\end{proof}

\subsection{The common-shock verdict (main-text Proposition 6)}\label{si:csproofs}

The variant is defined in main-text Section~4.6 and S4 below. Write $A(p)=ps_1+(1-p)s_2$ for the aggregate margin at weight share $p$.

\begin{proof}[Proof sketch (machine-anchored)]
(i) \act{fast}'s common-class requirement is $A(p)\ge\max(2p,\ p+\tfrac12)$ (its own tent $3/2$ in $s_1$ plus the $s_1$-event, and the $s_2$-event weighted by $1-p$); \act{slow}'s is the mirror; \act{staged}'s ($x\ge1$) is $A(p)\ge\max(p,1-p)/2-1/8$, the event exposure of a growing floor (kinked in $p$; its admissibility decomposes exactly into the two affine conditions of the kink's branches). The pointwise minimum of the two principal requirements is a tent with its only kink at $p=\tfrac12$, so an affine $A$ covers it exactly at the three abscissae $p\in\{0,\tfrac12,1\}$; the closed forms follow: typed acceptance needs $s_1\ge2$ with $s_2\ge\tfrac12$ (or the mirror) or the staged branch $x\ge1$, $\min(s)\ge\tfrac38$; weak acceptance with $x<1$ needs $s_1\ge\tfrac12$, $s_2\ge\tfrac12$ (the endpoints $p=0,1$) and $s_1+s_2\ge2$ (the kink). The licensing endpoint of \act{fast} on the surviving region equals its action-indexed endpoint $1/(1+\rho_1)$, since the binding branch there is the own-tent branch. The verifier checks the closed forms against independent grid computations ($5{,}000$-state grids) plus a curated witness list, both directions.
(ii) The augmented menu adds $\act{staged}_\kappa$, admissible at $x+\kappa\ge1$ with the same floor requirement $s_i\ge\tfrac38$; $\kappa$ cannot move the floors, so on $\{\min(s)<\tfrac38\}\setminus\Vtyp^{\mathrm{cs}}$ no finite $\kappa$ certifies. On the complement the requirement is the action-indexed one, $\kappa=(1-x)_+$. The $3/8$ boundary is structural: the event exposure of a growing floor is $\max(0,\delta-1/8)=3/8$ at $\delta=1/2$.
(iii) The blend's $s_1$-trough is $s_1-\delta-\delta_{\mathrm{blend}}\cdot\tfrac32$ and its $s_2$-trough $s_2-\delta-(1-\delta_{\mathrm{blend}})\cdot\tfrac32$ (each coordinate takes its event plus the corresponding fraction of the principal tent); the window is the nonnegativity interval of the two intersected with $[0,1]$, nonempty iff $s_1\ge\delta$, $s_2\ge\delta$ and $s_1+s_2\ge\tfrac32+2\delta$.
\end{proof}

